\subsection{Danh sách đặc tả Use-case}

Phần này đặc tả chi tiết các use-case chính ở mức nghiệp vụ của toàn bộ hệ thống. Mỗi use-case được trình bày dưới dạng bảng theo đúng cấu trúc của báo cáo mẫu để thuận tiện cho việc thiết kế API, giao diện và kiểm thử.

\usecasespec{UC-01}{Đăng ký tài khoản}{Khách}{Cho phép khách tạo mới tài khoản Mingo bằng email hoặc số điện thoại sau khi hoàn tất xác thực OTP trước đăng ký.}{Định danh đăng ký chưa tồn tại trong hệ thống; dịch vụ OTP hoạt động bình thường; người dùng chưa có phiên đăng nhập bắt buộc.}{Tài khoản mới được tạo, dữ liệu hồ sơ ban đầu được khởi tạo và người dùng có thể tiếp tục đăng nhập.}{\begin{minipage}[t]{0.98\linewidth}\raggedright 1. Khách nhập tên hiển thị, email hoặc số điện thoại, mật khẩu và các trường bắt buộc.\par 2. Hệ thống tiếp nhận yêu cầu gửi OTP trước đăng ký.\par 3. Dịch vụ OTP hoặc Email hoặc SMS phát hành và gửi mã xác thực.\par 4. Khách nhập mã OTP đã nhận.\par 5. Hệ thống kiểm tra mã, thời hạn hiệu lực và giới hạn thử lại.\par 6. Hệ thống tạo tài khoản mới và lưu thông tin đăng ký.\par 7. Hệ thống trả kết quả đăng ký thành công cho client.\end{minipage}}{\begin{minipage}[t]{0.98\linewidth}\raggedright 1. Email hoặc số điện thoại đã tồn tại, hệ thống từ chối đăng ký.\par 2. OTP sai, hết hạn hoặc vượt giới hạn thử lại, hệ thống yêu cầu gửi mã mới.\par 3. Dịch vụ gửi OTP lỗi, hệ thống trả thông báo thất bại và không tạo tài khoản.\end{minipage}}

\usecasespec{UC-02}{Đăng nhập và duy trì phiên}{Khách, Người dùng}{Cho phép người dùng đăng nhập bằng tài khoản cục bộ hoặc Google, hoàn tất 2FA nếu được bật và duy trì phiên làm việc bằng refresh token.}{Tài khoản tồn tại, chưa bị khóa, còn hiệu lực và có phương thức đăng nhập hợp lệ.}{Người dùng có phiên đăng nhập hợp lệ hoặc phiên hiện tại bị thu hồi khi đăng xuất.}{\begin{minipage}[t]{0.98\linewidth}\raggedright 1. Người dùng nhập email hoặc số điện thoại cùng mật khẩu, hoặc chọn đăng nhập Google.\par 2. Hệ thống xác thực thông tin đầu vào.\par 3. Nếu tài khoản chưa bật 2FA, hệ thống phát hành access token và refresh token.\par 4. Nếu tài khoản đã bật 2FA, hệ thống yêu cầu mã xác thực bổ sung.\par 5. Người dùng nhập mã 2FA.\par 6. Hệ thống kiểm tra mã và phát hành bộ token nếu hợp lệ.\par 7. Trong quá trình sử dụng, client có thể gọi refresh token để duy trì phiên.\par 8. Khi người dùng đăng xuất, hệ thống thu hồi hoặc xóa refresh token hiện hành.\end{minipage}}{\begin{minipage}[t]{0.98\linewidth}\raggedright 1. Sai mật khẩu hoặc tài khoản không tồn tại, hệ thống từ chối đăng nhập.\par 2. Mã 2FA sai hoặc hết hạn, hệ thống không cấp phiên mới.\par 3. Refresh token không tồn tại hoặc không hợp lệ, hệ thống yêu cầu đăng nhập lại.\end{minipage}}

\usecasespec{UC-03}{Quản lý xác thực đa yếu tố}{Người dùng}{Cho phép người dùng thiết lập, kích hoạt hoặc vô hiệu hóa xác thực hai yếu tố cho tài khoản của mình.}{Người dùng đã đăng nhập và có phiên xác thực hợp lệ.}{Trạng thái 2FA của tài khoản được cập nhật đúng theo thao tác của người dùng.}{\begin{minipage}[t]{0.98\linewidth}\raggedright 1. Người dùng mở màn hình bảo mật tài khoản.\par 2. Hệ thống sinh secret và dữ liệu QR để cấu hình 2FA.\par 3. Người dùng quét mã QR hoặc nhập secret vào ứng dụng xác thực.\par 4. Người dùng nhập mã OTP sinh ra từ ứng dụng.\par 5. Hệ thống kiểm tra mã và bật 2FA cho tài khoản.\par 6. Khi cần, người dùng gửi yêu cầu tắt 2FA.\par 7. Hệ thống kiểm tra lại điều kiện an toàn và cập nhật trạng thái mới.\end{minipage}}{\begin{minipage}[t]{0.98\linewidth}\raggedright 1. Mã kích hoạt sai, hệ thống không bật 2FA.\par 2. Phiên đăng nhập hết hạn, hệ thống yêu cầu đăng nhập lại.\par 3. Điều kiện xác minh khi tắt 2FA không hợp lệ, hệ thống từ chối thao tác.\end{minipage}}

\usecasespec{UC-04}{Khôi phục mật khẩu}{Khách}{Hỗ trợ người dùng đặt lại mật khẩu khi quên thông tin đăng nhập bằng email hoặc số điện thoại đã liên kết.}{Tài khoản tồn tại và có kênh nhận mã xác minh hợp lệ.}{Mật khẩu mới được cập nhật, mật khẩu cũ mất hiệu lực và người dùng có thể đăng nhập lại.}{\begin{minipage}[t]{0.98\linewidth}\raggedright 1. Khách chọn chức năng quên mật khẩu.\par 2. Hệ thống yêu cầu nhập email hoặc số điện thoại.\par 3. Hệ thống tạo mã xác minh hoặc token đặt lại mật khẩu.\par 4. Dịch vụ OTP hoặc Email hoặc SMS gửi thông tin khôi phục.\par 5. Khách nhập mã xác minh và mật khẩu mới.\par 6. Hệ thống kiểm tra mã xác minh và chính sách mật khẩu.\par 7. Hệ thống cập nhật mật khẩu mới cho tài khoản.\end{minipage}}{\begin{minipage}[t]{0.98\linewidth}\raggedright 1. Mã xác minh sai hoặc hết hạn, hệ thống từ chối đổi mật khẩu.\par 2. Định danh không tồn tại, hệ thống trả thông báo tổng quát để tránh lộ dữ liệu tài khoản.\par 3. Mật khẩu mới không đạt yêu cầu độ mạnh, hệ thống yêu cầu nhập lại.\end{minipage}}

\usecasespec{UC-05}{Quản lý hồ sơ cá nhân}{Người dùng}{Cho phép người dùng xem, cập nhật và cá nhân hóa hồ sơ cá nhân trên nền tảng.}{Người dùng đã đăng nhập và tài khoản còn hoạt động.}{Thông tin hồ sơ, ảnh đại diện hoặc ảnh nền được cập nhật thành công trong hệ thống.}{\begin{minipage}[t]{0.98\linewidth}\raggedright 1. Người dùng mở trang hồ sơ cá nhân.\par 2. Hệ thống trả về thông tin hồ sơ hiện tại.\par 3. Người dùng chỉnh sửa tên hiển thị, bio và các trường cho phép.\par 4. Người dùng có thể tải ảnh đại diện hoặc ảnh nền mới.\par 5. Hệ thống tải tệp lên Cloudinary và nhận URL lưu trữ.\par 6. Hệ thống cập nhật hồ sơ người dùng và trả dữ liệu mới cho client.\end{minipage}}{\begin{minipage}[t]{0.98\linewidth}\raggedright 1. Tệp tải lên không hợp lệ hoặc vượt giới hạn, hệ thống từ chối cập nhật ảnh.\par 2. Phiên đăng nhập hết hạn, hệ thống yêu cầu người dùng đăng nhập lại.\end{minipage}}

\usecasespec{UC-06}{Tìm kiếm toàn cục}{Người dùng}{Cho phép người dùng tìm kiếm người dùng, bài viết và hashtag trong hệ thống để khám phá nội dung.}{Người dùng đã đăng nhập và có quyền truy cập khu vực tìm kiếm.}{Danh sách kết quả phù hợp được trả về cho client theo loại thực thể tương ứng.}{\begin{minipage}[t]{0.98\linewidth}\raggedright 1. Người dùng nhập từ khóa tìm kiếm.\par 2. Hệ thống nhận truy vấn và chuẩn hóa chuỗi tìm kiếm.\par 3. Hệ thống tìm kiếm đồng thời trên người dùng, bài viết và hashtag.\par 4. Hệ thống tổng hợp, sắp xếp và trả về kết quả cho client.\par 5. Người dùng chọn một kết quả để xem chi tiết hoặc điều hướng đến màn hình liên quan.\end{minipage}}{\begin{minipage}[t]{0.98\linewidth}\raggedright 1. Từ khóa rỗng hoặc không hợp lệ, hệ thống yêu cầu nhập lại.\par 2. Không có kết quả phù hợp, hệ thống trả danh sách rỗng.\end{minipage}}

\usecasespec{UC-07}{Quản lý quan hệ người dùng}{Người dùng}{Hỗ trợ người dùng xây dựng và kiểm soát quan hệ theo dõi, close friend và danh sách chặn.}{Người dùng đã đăng nhập và tài khoản đích tồn tại.}{Trạng thái quan hệ giữa hai người dùng được cập nhật chính xác và các bộ đếm liên quan được đồng bộ.}{\begin{minipage}[t]{0.98\linewidth}\raggedright 1. Người dùng gửi yêu cầu theo dõi tới một tài khoản khác.\par 2. Nếu tài khoản đích ở chế độ riêng tư, hệ thống tạo yêu cầu chờ phản hồi.\par 3. Tài khoản đích chấp nhận hoặc từ chối yêu cầu theo dõi.\par 4. Khi quan hệ theo dõi hình thành, người dùng có thể bỏ theo dõi hoặc xóa follower.\par 5. Người dùng có thể gửi yêu cầu close friend cho một quan hệ phù hợp.\par 6. Người dùng có thể chặn hoặc bỏ chặn tài khoản khác khi cần.\par 7. Hệ thống cập nhật thống kê follower, following và danh sách liên quan.\end{minipage}}{\begin{minipage}[t]{0.98\linewidth}\raggedright 1. Người dùng gửi yêu cầu trùng lặp, hệ thống từ chối thao tác.\par 2. Tài khoản đang nằm trong danh sách chặn, hệ thống không cho tạo quan hệ mới.\par 3. Người dùng không đủ quyền chấp nhận hoặc từ chối yêu cầu không thuộc về mình.\end{minipage}}

\usecasespec{UC-08}{Quản lý bài viết}{Người dùng}{Cho phép người dùng tạo, xem, chỉnh sửa và xóa bài viết cùng nội dung đa phương tiện.}{Người dùng đã đăng nhập và tài khoản còn hoạt động.}{Bài viết được tạo mới, cập nhật hoặc xóa; các pipeline nền được kích hoạt khi cần.}{\begin{minipage}[t]{0.98\linewidth}\raggedright 1. Người dùng nhập nội dung bài viết, hashtag, mention, quyền riêng tư và media đính kèm nếu có.\par 2. Hệ thống tải media lên Cloudinary nếu bài viết có ảnh hoặc video.\par 3. Hệ thống lưu bài viết vào cơ sở dữ liệu.\par 4. Sau khi lưu thành công, hệ thống khởi chạy pipeline kiểm duyệt nội dung và phân tích thuật ngữ.\par 5. Người dùng có thể xem danh sách bài viết, chi tiết bài viết hoặc bài viết theo hồ sơ người dùng.\par 6. Tác giả có thể chỉnh sửa hoặc xóa bài viết của mình.\end{minipage}}{\begin{minipage}[t]{0.98\linewidth}\raggedright 1. Media tải lên lỗi, hệ thống từ chối tạo bài hoặc yêu cầu tải lại tệp.\par 2. Người dùng sửa hoặc xóa bài viết không thuộc quyền sở hữu, hệ thống từ chối thao tác.\par 3. Bài viết bị đánh dấu vi phạm nghiêm trọng, nội dung có thể bị ẩn khỏi feed.\end{minipage}}

\usecasespec{UC-09}{Tương tác với nội dung}{Người dùng}{Cho phép người dùng thể hiện phản hồi với nội dung trong hệ thống và tạo tín hiệu hành vi phục vụ cá nhân hóa.}{Người dùng đã đăng nhập; bài viết hoặc bình luận đích tồn tại và có quyền tương tác.}{Dữ liệu tương tác và tín hiệu hành vi được ghi nhận, đồng thời có thể phát sinh thông báo cho chủ sở hữu nội dung.}{\begin{minipage}[t]{0.98\linewidth}\raggedright 1. Người dùng mở một bài viết hoặc bình luận.\par 2. Người dùng thực hiện thao tác thích, bỏ thích, lưu, bỏ lưu, bình luận, trả lời, chia sẻ, gửi qua tin nhắn hoặc repost.\par 3. Hệ thống cập nhật dữ liệu tương tác tương ứng.\par 4. Nếu thao tác ảnh hưởng đến gợi ý nội dung, hệ thống ghi nhận thêm sự kiện hành vi.\par 5. Hệ thống trả trạng thái mới và có thể phát sinh thông báo cho chủ sở hữu nội dung.\end{minipage}}{\begin{minipage}[t]{0.98\linewidth}\raggedright 1. Nội dung không tồn tại hoặc đã bị xóa, hệ thống từ chối thao tác.\par 2. Người dùng gửi bình luận rỗng hoặc nội dung không hợp lệ, hệ thống yêu cầu chỉnh sửa.\par 3. Chia sẻ qua tin nhắn thất bại do hộp thoại đích không hợp lệ.\end{minipage}}

\usecasespec{UC-10}{Xem feed cá nhân hóa}{Người dùng}{Cung cấp feed Friends và Explore phù hợp với sở thích và hành vi của người dùng.}{Người dùng đã đăng nhập; hệ thống có dữ liệu bài viết hợp lệ để xếp hạng.}{Danh sách bài viết được trả về cùng thông tin phân trang và các chỉ số liên quan.}{\begin{minipage}[t]{0.98\linewidth}\raggedright 1. Người dùng mở màn hình feed và chọn tab Friends hoặc Explore.\par 2. Hệ thống xác định nguồn dữ liệu ứng với tab đã chọn.\par 3. Hệ thống tính điểm và sắp xếp bài viết dựa trên hồ sơ hành vi, chủ đề, độ phổ biến và tín hiệu thời gian.\par 4. Hệ thống trả danh sách bài viết đã phân trang cho client.\par 5. Người dùng có thể gửi phản hồi như see more, hide hoặc not interested.\par 6. Hệ thống ghi nhận phản hồi để cập nhật hồ sơ gợi ý ở các lần truy xuất sau.\end{minipage}}{\begin{minipage}[t]{0.98\linewidth}\raggedright 1. Hồ sơ hành vi chưa đủ dữ liệu, hệ thống fallback sang tín hiệu mặc định như độ mới và độ phổ biến.\par 2. Không có bài viết phù hợp, hệ thống trả danh sách rỗng.\end{minipage}}

\usecasespec{UC-11}{Giải thích thuật ngữ văn hóa mạng}{Người dùng, Gemini API}{Hỗ trợ người dùng hiểu nội dung có chứa slang hoặc tham chiếu văn hóa mạng theo đúng ngữ cảnh bài viết.}{Bài viết đã được lưu hoặc đã đi qua pipeline phân tích thuật ngữ.}{Thuật ngữ được phát hiện, gắn vị trí trong nội dung và có giải thích để hiển thị ở client.}{\begin{minipage}[t]{0.98\linewidth}\raggedright 1. Sau khi bài viết được tạo hoặc chỉnh sửa, hệ thống dò tìm thuật ngữ theo từ điển slang hiện hành.\par 2. Hệ thống gom các thuật ngữ phát hiện được và gọi Gemini API để lấy giải thích theo ngữ cảnh.\par 3. Hệ thống ghép kết quả giải thích với vị trí xuất hiện trong nội dung và lưu vào bài viết.\par 4. Khi người dùng xem bài viết, client tải dữ liệu cultural terms từ backend.\par 5. Người dùng nhấn vào thuật ngữ được highlight để xem nghĩa, nguồn gốc và ghi chú ngữ cảnh.\par 6. Nếu phát hiện giải thích sai hoặc không phù hợp, người dùng có thể gửi báo lỗi thuật ngữ.\end{minipage}}{\begin{minipage}[t]{0.98\linewidth}\raggedright 1. Không phát hiện thuật ngữ nào, hệ thống kết thúc pipeline mà không gọi AI.\par 2. Gemini API lỗi hoặc trả dữ liệu sai định dạng, hệ thống fallback về kết quả rỗng để không làm hỏng luồng tạo bài.\par 3. Người dùng yêu cầu phân tích lại, hệ thống xóa kết quả cũ và chạy pipeline mới.\end{minipage}}

\usecasespec{UC-12}{Nhận và quản lý thông báo}{Người dùng, Pusher/Socket.IO}{Cho phép người dùng theo dõi các sự kiện liên quan đến tài khoản và nội dung của mình.}{Người dùng đã đăng nhập; thiết bị đã được đăng ký nếu dùng push notification.}{Thiết bị được đăng ký hoặc gỡ đăng ký; trạng thái thông báo được cập nhật đúng.}{\begin{minipage}[t]{0.98\linewidth}\raggedright 1. Người dùng đăng ký token thiết bị để nhận thông báo đẩy.\par 2. Khi có sự kiện như tương tác, theo dõi, nhắn tin hoặc xử lý báo cáo, hệ thống tạo thông báo tương ứng.\par 3. Người dùng mở danh sách thông báo.\par 4. Hệ thống trả danh sách thông báo, số lượng chưa đọc và trạng thái đã xem.\par 5. Người dùng đánh dấu đã đọc, đánh dấu đã xem hoặc xóa thông báo.\par 6. Hệ thống cập nhật trạng thái và đồng bộ dữ liệu với client.\end{minipage}}{\begin{minipage}[t]{0.98\linewidth}\raggedright 1. Token thiết bị không hợp lệ, hệ thống từ chối đăng ký.\par 2. Người dùng thao tác trên thông báo không thuộc về mình, hệ thống từ chối.\end{minipage}}

\usecasespec{UC-13}{Nhắn tin và gọi thoại}{Người dùng, Cloudinary, Pusher/Socket.IO}{Cung cấp kênh giao tiếp thời gian thực giữa các người dùng dưới dạng hội thoại trực tiếp hoặc nhóm.}{Người dùng đã đăng nhập; hội thoại trực tiếp hoặc nhóm chat tồn tại hoặc có thể tạo mới.}{Tin nhắn, trạng thái đọc, trạng thái trực tuyến hoặc lịch sử cuộc gọi được cập nhật.}{\begin{minipage}[t]{0.98\linewidth}\raggedright 1. Người dùng mở danh sách cuộc trò chuyện hoặc tạo nhóm chat mới.\par 2. Hệ thống trả danh sách direct box, group box hoặc chi tiết nhóm.\par 3. Người dùng gửi tin nhắn văn bản hoặc tệp đính kèm.\par 4. Nếu có tệp, hệ thống tải lên nơi lưu trữ phù hợp và tạo bản ghi tin nhắn.\par 5. Hệ thống phát tin nhắn mới qua kênh thời gian thực cho các thành viên liên quan.\par 6. Người dùng có thể sửa, thu hồi hoặc xóa tin nhắn theo quyền.\par 7. Người dùng có thể đánh dấu đã đọc, tìm kiếm trong hộp thoại, xem media đã gửi, cập nhật thông tin nhóm, thêm hoặc xóa thành viên, rời nhóm hoặc phân quyền quản trị nhóm.\par 8. Người dùng có thể khởi tạo cuộc gọi và cập nhật trạng thái cuộc gọi.\end{minipage}}{\begin{minipage}[t]{0.98\linewidth}\raggedright 1. Người dùng không thuộc hộp thoại hoặc không có quyền quản trị nhóm, hệ thống từ chối thao tác quản lý nhóm.\par 2. Tệp vượt giới hạn kích thước hoặc không tải lên được, hệ thống từ chối gửi tin nhắn đính kèm.\par 3. Tin nhắn đã bị thu hồi hoặc xóa, hệ thống không cho chỉnh sửa tiếp.\end{minipage}}

\usecasespec{UC-14}{Gửi báo cáo vi phạm}{Người dùng}{Cho phép người dùng chủ động phản ánh nội dung hoặc tài khoản có dấu hiệu vi phạm.}{Người dùng đã đăng nhập; đối tượng bị báo cáo tồn tại trong hệ thống.}{Báo cáo được lưu và sẵn sàng cho pipeline xử lý hoặc xem xét thủ công.}{\begin{minipage}[t]{0.98\linewidth}\raggedright 1. Người dùng chọn báo cáo bài viết, bình luận, tài khoản hoặc thuật ngữ.\par 2. Hệ thống yêu cầu chọn lý do và nhập mô tả bổ sung nếu cần.\par 3. Người dùng xác nhận gửi báo cáo.\par 4. Hệ thống lưu báo cáo, liên kết với đối tượng liên quan và người gửi báo cáo.\par 5. Tùy loại báo cáo, hệ thống tăng chỉ số report hoặc đưa đối tượng vào hàng đợi xem xét.\par 6. Hệ thống trả kết quả gửi báo cáo thành công.\end{minipage}}{\begin{minipage}[t]{0.98\linewidth}\raggedright 1. Người dùng báo cáo trùng lặp trong điều kiện không cho phép, hệ thống từ chối hoặc gộp báo cáo.\par 2. Đối tượng không tồn tại, hệ thống từ chối tạo báo cáo.\end{minipage}}

\usecasespec{UC-15}{Kiểm duyệt nội dung bằng AI}{Hệ thống, Gemini API, Quản trị viên}{Phát hiện nội dung độc hại, spam hoặc thù ghét để giảm rủi ro trên nền tảng bằng pipeline hai tầng.}{Bài viết hoặc bình luận đã được tạo và sẵn sàng cho pipeline kiểm duyệt.}{Nội dung được phân loại thành duyệt ngay, gắn cờ cần xem xét hoặc tự động ẩn.}{\begin{minipage}[t]{0.98\linewidth}\raggedright 1. Khi người dùng tạo hoặc cập nhật nội dung, hệ thống khởi chạy tầng kiểm tra rule-based.\par 2. Hệ thống kiểm tra blocklist, mẫu hate speech và tín hiệu spam.\par 3. Nếu phát hiện vi phạm rõ ràng, hệ thống gắn trạng thái vi phạm ngay.\par 4. Nếu nội dung thuộc vùng nghi ngờ hoặc thỏa điều kiện cần AI, hệ thống gọi Gemini API.\par 5. Hệ thống nhận điểm toxic, hate speech và spam từ AI, sau đó kết hợp với tín hiệu rule-based.\par 6. Hệ thống ra quyết định cuối cùng: duyệt, gắn cờ human review hoặc tự động ẩn.\par 7. Quản trị viên có thể tiếp tục xem xét các nội dung bị gắn cờ hoặc bị báo cáo.\end{minipage}}{\begin{minipage}[t]{0.98\linewidth}\raggedright 1. Gemini API lỗi hoặc timeout, hệ thống dùng cơ chế fallback để không làm thất bại thao tác tạo nội dung.\par 2. Nội dung bị báo cáo nhiều lần, hệ thống có thể tăng mức ưu tiên kiểm duyệt.\par 3. Kết quả AI không hợp lệ, hệ thống dùng điểm fallback thấp và ghi nhận lý do.\end{minipage}}

\usecasespec{UC-16}{Quản trị báo cáo và nội dung}{Quản trị viên}{Hỗ trợ quản trị viên kiểm soát rủi ro nội dung và người dùng trên nền tảng.}{Quản trị viên đã đăng nhập với quyền hợp lệ.}{Báo cáo được xử lý; bài viết, bình luận hoặc tài khoản được cập nhật trạng thái quản trị tương ứng.}{\begin{minipage}[t]{0.98\linewidth}\raggedright 1. Quản trị viên mở dashboard quản trị hoặc hàng đợi báo cáo.\par 2. Hệ thống trả danh sách báo cáo đang chờ, bài viết, bình luận và người dùng liên quan.\par 3. Quản trị viên xem chi tiết đối tượng bị báo cáo và lịch sử hoạt động liên quan.\par 4. Quản trị viên chọn hành động như ẩn bài, hiện lại bài, xóa bài, xử lý bình luận, khóa hoặc mở khóa người dùng, đánh dấu báo cáo đã giải quyết.\par 5. Hệ thống lưu kết quả xử lý và ghi nhận audit log.\par 6. Hệ thống trả trạng thái cập nhật cho trang quản trị.\end{minipage}}{\begin{minipage}[t]{0.98\linewidth}\raggedright 1. Quản trị viên không có quyền trên thao tác cụ thể, hệ thống từ chối.\par 2. Đối tượng đã bị xóa hoặc xử lý trước đó, hệ thống trả trạng thái hiện tại để tránh thao tác lặp.\end{minipage}}

\usecasespec{UC-17}{Quản lý từ điển slang}{Quản trị viên}{Duy trì danh mục thuật ngữ văn hóa mạng phục vụ phát hiện và giải thích trong bài viết.}{Quản trị viên đã đăng nhập và có quyền quản lý từ điển hệ thống.}{Từ điển slang được cập nhật, thêm mới hoặc bật tắt đúng mục cần quản lý.}{\begin{minipage}[t]{0.98\linewidth}\raggedright 1. Quản trị viên mở màn hình quản lý từ điển slang.\par 2. Hệ thống trả danh sách các mục slang hiện có cùng trạng thái hoạt động.\par 3. Quản trị viên thêm một mục mới gồm thuật ngữ, regex pattern, nghĩa và thông tin mô tả.\par 4. Hệ thống kiểm tra dữ liệu đầu vào và lưu mục mới.\par 5. Quản trị viên có thể bật hoặc tắt một mục để điều chỉnh phạm vi phát hiện.\par 6. Hệ thống cập nhật từ điển và làm mới cache theo chu kỳ hệ thống.\end{minipage}}{\begin{minipage}[t]{0.98\linewidth}\raggedright 1. Regex hoặc dữ liệu mục slang không hợp lệ, hệ thống từ chối lưu.\par 2. Quản trị viên thao tác trên mục không tồn tại, hệ thống trả lỗi phù hợp.\end{minipage}}

\usecasespec{UC-18}{Theo dõi thống kê hệ thống}{Quản trị viên}{Cung cấp cho quản trị viên cái nhìn tổng quan về trạng thái vận hành và hiệu quả của hệ thống.}{Quản trị viên đã đăng nhập và có quyền truy cập dashboard quản trị.}{Quản trị viên nhận được các chỉ số và dữ liệu giám sát cần thiết để theo dõi và ra quyết định.}{\begin{minipage}[t]{0.98\linewidth}\raggedright 1. Quản trị viên truy cập dashboard hoặc khu vực thống kê.\par 2. Hệ thống trả các số liệu tổng quan về người dùng, bài viết, vi phạm, báo cáo và tương tác.\par 3. Quản trị viên xem chỉ số hiệu năng AI, thống kê vi phạm theo ngày và audit log.\par 4. Khi cần, quản trị viên mở sâu vào thống kê của từng người dùng, bài viết hoặc nhóm dữ liệu liên quan.\end{minipage}}{\begin{minipage}[t]{0.98\linewidth}\raggedright 1. Dữ liệu thống kê tạm thời chưa sẵn sàng, hệ thống trả kết quả rỗng hoặc giá trị mặc định.\par 2. Người dùng không có quyền quản trị cố truy cập chức năng này, hệ thống từ chối.\end{minipage}}

